\newcounter{defctr}[section]
\renewcommand{\thedefctr}{\thesection.\arabic{defctr}}

\section{Cosmological Operational Closure and the Scale $\ell = c/H$}
  \label{app:cosmo_resolution}

  This appendix derives Eq.~\eqref{eq:ell_cosmo} used in Section~\ref{subsec:cosmo_resolution}.
  The derivation relies solely on the operational closure axiom and the
  homogeneous saturated kinematics of the projected state.
  No geometric horizon concept is invoked.

  \subsection{A.1 Operational closure and descriptive resolution}
    \label{subsec:operational-closure}

    \refstepcounter{defctr}\label{def:ell_closure}
    \noindent\textbf{Definition~\thedefctr\ (Operational closure and descriptive resolution).}
    The descriptive resolution $\ell$ is the minimal coarse-graining neighbourhood such that the update of the
    projected state $U$ at epoch $t$ depends on all substrate channels that can influence the update within a
    single admissible step, and on no others.
    \par\medskip

    \noindent By definition~\ref{def:ell_closure}, the descriptive resolution $\ell$
    is the minimal coarse-graining neighbourhood such that the update of the
    projected state $U$ at epoch $t$ depends on all substrate channels that can
    influence the update within a single admissible step, and on no others.

    Closure therefore requires two properties:

    \begin{enumerate}
      \item Completeness: all admissible influences must be included.
      \item Minimality: no inadmissible or causally irrelevant regions are included.
    \end{enumerate}

  \subsection{Homogeneous saturated kinematics}
    \label{subsec:homogeneous-saturated-kinematics}

    In the homogeneous late-time regime,
    \begin{equation}
      D_{\mathrm{loc}}\chi_{\mathrm{eff}} = c,
    \end{equation}
    so propagation of relational influence is bounded by $c$.

    The projected state evolves according to
    \begin{equation}
      H(t) = \frac{\dot{\chi}}{\chi}.
    \end{equation}

    In a homogeneous regime, $H(t)$ is the unique intrinsic dimensionless rate.
    The characteristic update time is therefore
    \begin{equation}
      \tau(t) \sim \left(\frac{\dot{\chi}}{\chi}\right)^{-1} = H(t)^{-1}.
    \end{equation}

    This is not imposed externally.
    It follows from the requirement that an update corresponds to an
    order-one change of the dimensionless global state variable.

  \subsection{Influence domain and lower bound}
    \label{subsec:influence-domain-and-lower-bound}

    Within one update step of duration $\tau(t)$, the maximal
    propagation distance of relational influence is
    \begin{equation}
      r_{\mathrm{max}} = c \tau(t) = \frac{c}{H(t)}.
    \end{equation}

    Any substrate configuration at distance $r \le r_{\mathrm{max}}$
    can influence the projected state within the update step.
    Excluding such configurations from the coarse-graining neighbourhood
    would violate completeness of closure.

    Therefore,
    \begin{equation}
      \ell \ge \frac{c}{H(t)}.
    \end{equation}

  \subsection{Minimality and upper bound}
    \label{subsec:minimality-and-upper-bound}

    Conversely, any configuration at distance $r > r_{\mathrm{max}}$
    cannot influence the update within the admissible step.

    Including such regions in the neighbourhood does not modify the update
    map of $U$, but enlarges the coarse-graining domain beyond the minimal
    operationally closed set.

    Minimality of closure therefore requires
    \begin{equation}
      \ell \le \frac{c}{H(t)}.
    \end{equation}

  \subsection{Uniqueness of the homogeneous resolution}
    \label{subsec:uniqueness-of-the-homogeneous-resolution}

    Combining the bounds yields
    \begin{equation}
      \ell_{\mathrm{cosmo}}(t) = \frac{c}{H(t)}.
    \end{equation}

    No additional scale is available in the homogeneous regime.
    Any alternative choice $\ell = f(H)$ would either violate dimensional
    consistency or introduce an extraneous parameter, contradicting
    the ontological minimality of the framework.

  \subsection{Consequence for the saturation scale}
    \label{subsec:consequence-for-the-saturation-scale}

    Substituting into
    \begin{equation}
      a_\star(t) \sim \kappa \frac{c_{\mathrm{BI}}^2}{\ell(t)}
    \end{equation}
    gives
    \begin{equation}
      a_\star(t) \sim \kappa\, \frac{c_{\mathrm{BI}}^2}{c}\, H(t).
    \end{equation}

    In projectable regimes $c_{\mathrm{BI}} \to c$, yielding
    \begin{equation}
      a_\star(t) \sim \kappa\, c H(t).
    \end{equation}

    This establishes a structural link between the galactic
    saturation scale and the cosmological expansion rate.
    The normalization factor $\kappa$ is fixed independently by the
    $\mathrm{U}(1)$ holonomy condition discussed in Section~\ref{subsec:holonomy_normalization}.
